\scanpage{33}
% @ 1
E\uline{xemples de groupes algébriques affines sur} $k$.

% @ 2–8
\diagram{\scriptsize\(\begin{array}{@{}l@{\;}l@{\;}l@{\;}c@{\;}l@{}}
G_m, & G_a, & \mathrm{GL}(n), & \mathrm{O}(Q)\unsure{.}, &
 \text{ss-groupe de }\mathrm{GL}(n)\ \text{invariant des tenseurs}\\[-.5ex]
&&\mathrm{SL}(n)&\begin{array}{c}\cap\\[-.8ex]\mathrm{GL}(n)\end{array}
 &\hfill .\,.\,.\,.\,. \\
&&\mathrm{Tr}(n)&&\\
\unsure{\mathrm{Sp}.|\mathrm{Ex}.}&&\mathrm{Tr}_{\unsure{0|o}}(n)&&\\
\multicolumn{2}{r}{\mathrm{Diag}_0(n)} & \multicolumn{1}{l}{\subset\mathrm{Diag}(n)} &
\multicolumn{2}{l}{\text{immersion de }G_a\text{ dans }\mathbb{G}\unsure{\mathbb{L}|\mathrm{L}}(2)}\\
\multicolumn{2}{c}{\vcenter{\hbox{\tikz[baseline]\node[rotate=90,inner sep=0pt]{\(\unsure{\simeq|\sim}\)};}}} &&&\\
\multicolumn{2}{c}{G_m}&&&
\end{array}\)}

% @ 9
Soit $M$ un $k$-module.
% @ 10–11
\[
 W(M):\unsure{\{|\varepsilon}\ k'\longmapsto M\otimes_k \unsure{k|K}'
 \qquad\text{\parbox{.4\linewidth}{(en fait un $\unsure{\mathcal{O}|\mathscr{O}}_k$-Module, pas seulement un groupe)}}
\]

% @ 11–13
$\cdot$ Pour que ce soit représentable, il faut que $W(M)$ transforme injections $k'\unsure{-|\to} k''$ en injections, donc que $M$ soit p\uline{lat} sur $k$. \deleted{\unsure{So|Si}}

% @ 14–15
Il faut qu'il transforme produits (infinis) d'algèbres en produits
\[
 M\otimes_k\prod k'_i\longrightarrow\prod_i M\otimes_k k'_i
 \qquad\text{bijectif.}
\]

% @ 16–18
\deleted{On voit que ceci} faisant tous les $k'_i$ égaux \unsure{:|à} $k$, on trouve
\[
 M\otimes k^I\xrightarrow{\sim}\deleted{\unsure{\prod_i|\prod}} M^I
\]
% @ 19–21
Et on voit tout de suite que celà implique que $M$ est de type fini, \unsure{donc} \deleted{\unsure{un|une}} quotient d'un libre \deleted{\unsure{fini|f}} \deleted{\unsure{$R\mathrel{-}L$}} \deleted{\unsure{comme|librement}} on a $R\mathrel{-}L\mathrel{-}M\mathrel{-}0$
% @ 22–25
et comme $M$ est plat, \unsure{cel|celà} reste exact par $\otimes_k k^I$, donc $R\otimes k^I\xrightarrow{\sim}R^I$, donc $R$ de type fini, donc $M$ de présentation finie, donc $M$ plat de prés. finie, donc $M$ projectif de t-f. (facteur direct d'un libre)

% @ 26–27
Inversement, si $M$ projectif de t.f. \unsure{i.e.} facteur direct d'un libre $L$\inserted{$\simeq k^{\unsure{n}}$}, par projecteur $p$, on a
% @ 28–29
\[
 \begin{aligned}
 M\otimes_k k'&=\mathrm{Ker}\big((1-p)_{k'}:k'^{n}\to k'^{n}\big)\\
 W(M)&=\mathrm{Ker}\big(\mathbb{E}^{n}\overset{1-p}{\unsure{\relbar|\to}}\mathbb{E}^{n}\big)
 \end{aligned}
\]
% @ 30
c'est repr.

% archive sheet number: 32
